\subsection{Block Diagram Algebra --- Connecting System Components}

Block diagrams serve as the primary visual language for representing how individual system components interconnect to form complete control systems. While the transfer function approach introduced in Section 1.4.2 describes the input-output behavior of individual blocks, real control systems consist of multiple blocks connected in various configurations that must be analyzed as integrated wholes. Understanding how to manipulate and simplify block diagrams is therefore essential for analyzing any practical control system. The rules of block diagram algebra provide systematic procedures for combining individual transfer functions according to how their corresponding blocks are interconnected, allowing us to derive the overall transfer function of complex systems from their constituent parts.

A block diagram represents a dynamic system as a collection of functional blocks connected by signal lines. Each block contains the transfer function relating its input to its output, and the signal lines indicate the direction of information flow through the system. The power of block diagrams lies in their ability to visually represent the causal relationships between different parts of a system while simultaneously encoding the mathematical relationships that govern their behavior. When we draw a block diagram, we are creating both a visual model and a mathematical model simultaneously, with the understanding that we can analyze the system by applying algebraic rules to the transfer functions in accordance with the interconnection structure shown in the diagram.

Consider first the simplest type of interconnection: two blocks connected in series, also called cascade connection. In this configuration, the output of the first block becomes the input to the second block, so that the signal flows sequentially through both transfer functions. Suppose we have a block with transfer function $G_1(s)$ followed by a block with transfer function $G_2(s)$. If the input to the first block is $U(s)$ and its output is $Y_1(s)$, then $Y_1(s) = G_1(s)U(s)$. The second block takes $Y_1(s)$ as its input and produces $Y(s) = G_2(s)Y_1(s)$. Substituting the expression for $Y_1(s)$ into this equation yields $Y(s) = G_2(s)G_1(s)U(s)$. This result reveals that two blocks in series can be replaced by a single equivalent block whose transfer function is simply the product of the individual transfer functions, with the order of multiplication corresponding to the order of the blocks in the cascade. Importantly, the commutative property of multiplication ensures that $G_1(s)G_2(s) = G_2(s)G_1(s)$, so the equivalent transfer function is independent of the order in which we multiply the individual transfer functions. However, we must be cautious: this mathematical commutativity does not imply that the physical systems can be interchanged, as their internal states and energy storage elements may create dependencies that are not captured by transfer function representation alone.

The parallel connection represents a fundamentally different type of interconnection. In parallel configuration, the same input signal is applied simultaneously to multiple blocks, and their individual outputs are summed together to produce the system output. If two blocks with transfer functions $G_1(s)$ and $G_2(s)$ are connected in parallel, and both receive the same input $U(s)$, then the first block produces output $Y_1(s) = G_1(s)U(s)$ and the second produces $Y_2(s) = G_2(s)U(s)$. The system output is the sum of these individual outputs, so $Y(s) = Y_1(s) + Y_2(s) = \left[G_1(s) + G_2(s)\right]U(s)$. This derivation shows that parallel blocks combine through addition of their transfer functions, meaning the equivalent single block representing a parallel connection has transfer function $G_{\text{eq}}(s) = G_1(s) + G_2(s)$. The parallel configuration is commonly encountered when multiple control actions contribute to the same actuator, such as when proportional, integral, and derivative controllers are combined in a PID structure. The ability to combine such parallel paths into an equivalent transfer function greatly simplifies analysis of systems with multiple concurrent control paths.

The feedback interconnection represents the most important and commonly encountered configuration in control systems, as it forms the structural basis of virtually all closed-loop control systems. In a feedback configuration, a portion of the system output is fed back and subtracted from the reference input to produce an error signal that drives the forward-path blocks. Understanding how to derive the closed-loop transfer function for this configuration is essential for analyzing any feedback control system. Consider a system with forward-path transfer function $G(s)$ and feedback-path transfer function $H(s)$. The reference input $R(s)$ enters the system and is compared with the feedback signal $B(s)$ at a summing junction. The error signal $E(s)$ is the difference between reference and feedback: $E(s) = R(s) - B(s)$. The forward-path block produces output $Y(s) = G(s)E(s)$. The feedback signal is obtained by passing this output through the feedback transfer function: $B(s) = H(s)Y(s)$.

To derive the closed-loop transfer function relating output $Y(s)$ to reference input $R(s)$, we substitute the various relationships into each other. From the definitions above, we have $Y(s) = G(s)E(s)$ and $E(s) = R(s) - B(s)$. Substituting $B(s) = H(s)Y(s)$ into the error expression gives $E(s) = R(s) - H(s)Y(s)$. Now substitute this expression for $E(s)$ into the forward-path equation: $Y(s) = G(s)\left[R(s) - H(s)Y(s)\right]$. Expanding this yields $Y(s) = G(s)R(s) - G(s)H(s)Y(s)$. Gathering all terms involving $Y(s)$ on one side gives $Y(s) + G(s)H(s)Y(s) = G(s)R(s)$, which factors as $Y(s)\left[1 + G(s)H(s)\right] = G(s)R(s)$. Solving for the ratio $Y(s)/R(s)$ produces the closed-loop transfer function:

\begin{equation}
T(s) = \frac{Y(s)}{R(s)} = \frac{G(s)}{1 + G(s)H(s)}
\end{equation}

This result is the fundamental closed-loop transfer function for a unity feedback system when $H(s) = 1$, and it generalizes to the form shown when a non-unity feedback transfer function is present.

It is critically important to understand why a feedback interconnection cannot be reduced through simple multiplication of transfer functions. If we were to naively multiply the forward-path and feedback-path transfer functions, we would obtain $G(s)H(s)$, which represents the open-loop transfer function but not the closed-loop behavior. The distinction between open-loop and closed-loop transfer functions is not merely semantic but reflects fundamentally different system behaviors. The open-loop transfer function $G(s)H(s)$ describes the relationship between the feedback signal and the error signal when the loop is broken, but it tells us nothing about how the loop operates when closed. The closed-loop transfer function, in contrast, captures the actual input-output behavior of the system with feedback active, including all the complex interactions between forward and feedback paths that emerge when the loop is completed. The denominator term $1 + G(s)H(s)$ in the closed-loop transfer function represents the effect of feedback on system dynamics, and this term has no analogue in simple series or parallel combinations. When $G(s)H(s)$ is large in magnitude, the closed-loop transfer function approaches $1/H(s)$, meaning the system output becomes largely independent of the forward-path dynamics and is instead determined primarily by the feedback path. This property of feedback to make system behavior insensitive to variations in forward-path parameters is precisely what makes closed-loop control so powerful, and it cannot be understood through simple transfer function multiplication.

\begin{table}[htbp]
\centering
\begin{tabular}{|c|c|c|c|}
\hline
\rowcolor{UofSCGarnet}
\textcolor{white}{\textbf{Configuration}} & \textcolor{white}{\textbf{Diagram Structure}} & \textcolor{white}{\textbf{Equivalent Transfer Function}} & \textcolor{white}{\textbf{Formula}} \\
\hline
Series (Cascade) & Blocks connected end-to-end & Single equivalent block & $G_{\text{eq}}(s) = G_1(s) \cdot G_2(s)$ \\
\hline
Parallel & Same input, outputs summed & Single equivalent block & $G_{\text{eq}}(s) = G_1(s) + G_2(s)$ \\
\hline
Feedback & Output fed back and subtracted & Closed-loop transfer function & $T(s) = \dfrac{G(s)}{1 + G(s)H(s)}$ \\
\hline
\end{tabular}
\captionsetup{justification=centering}
\caption{Summary of block diagram interconnection rules showing how individual transfer functions combine in different configurations.}
\end{table}

To build intuition for block diagram reduction, let us work through a detailed example involving a feedback system with a controller, plant, and sensor. Consider a position control system where the reference position $R(s)$ is compared with the measured position $Y(s)$ to produce an error signal. The controller is a proportional-derivative (PD) controller with transfer function $G_c(s) = K_p + K_d s$, where $K_p$ is the proportional gain and $K_d$ is the derivative gain. The plant represents the mechanical system being controlled and has transfer function $G_p(s) = \dfrac{1}{ms^2 + b s}$, where $m$ is the mass and $b$ is the damping coefficient. The position sensor provides feedback with transfer function $H(s) = K_s$, representing a constant sensor gain. We wish to derive the closed-loop transfer function relating output position to reference position.

The block diagram shows the reference input entering a summing junction where it is compared with the sensor output. The error signal drives the PD controller, whose output becomes the input to the plant. The plant output is the physical position, which is measured by the sensor and fed back to the summing junction. To find the closed-loop transfer function, we identify the forward-path transfer function as $G(s) = G_c(s)G_p(s) = (K_p + K_d s) \cdot \dfrac{1}{ms^2 + b s}$ and the feedback-path transfer function as $H(s) = K_s$. Applying the feedback formula derived above, the closed-loop transfer function is:

\begin{equation}
T(s) = \frac{G(s)}{1 + G(s)H(s)} = \frac{(K_p + K_d s) \dfrac{1}{ms^2 + b s}}{1 + (K_p + K_d s) \dfrac{1}{ms^2 + b s} \cdot K_s}
\end{equation}

Simplifying this expression algebraically, we first combine the terms in the numerator and denominator:

\begin{equation}
T(s) = \frac{\dfrac{K_p + K_d s}{ms^2 + b s}}{1 + \dfrac{K_s (K_p + K_d s)}{ms^2 + b s}} = \frac{K_p + K_d s}{ms^2 + b s + K_s (K_p + K_d s)}
\end{equation}

Expanding the denominator by distributing the $K_s$ term gives:

\begin{equation}
T(s) = \frac{K_p + K_d s}{ms^2 + b s + K_s K_p + K_s K_d s} = \frac{K_p + K_d s}{ms^2 + (b + K_s K_d) s + K_s K_p}
\end{equation}

This final expression represents the complete closed-loop transfer function of the position control system. Several important observations emerge from this derivation. The denominator polynomial has been modified by the feedback terms: the constant term now includes $K_s K_p$ (proportional gain times sensor gain) instead of zero, and the $s$ coefficient now includes $K_s K_d$ (derivative gain times sensor gain) in addition to the original damping $b$. These modifications demonstrate how feedback reshapes the system dynamics. The effective damping has increased from $b$ to $b + K_s K_d$, meaning the derivative control action has enhanced the system's damping characteristics. The presence of the $K_s K_p$ term in the constant coefficient means the system now has a non-zero steady-state term that influences its type (the number of integrators in the forward path), which affects steady-state error behavior. Importantly, none of these feedback effects would be captured by simply multiplying the controller, plant, and sensor transfer functions; the structure of the feedback interconnection fundamentally changes the system behavior in ways that require the complete block diagram reduction procedure to reveal.

A more complex block diagram reduction demonstrates the power of systematic simplification for analyzing multi-loop systems. Consider a system with two feedback loops: an inner velocity feedback loop and an outer position feedback loop. The plant consists of an integrator (representing the conversion from velocity to position) followed by a first-order system (representing the dynamics from force to velocity). The inner loop contains a velocity sensor with gain $K_v$ that feeds back the velocity signal to be subtracted from the controller output. The outer loop uses a position sensor with gain $K_p$ to close the main feedback loop. To reduce this block diagram, we begin with the innermost loop. The forward-path transfer function for the inner loop is the product of the first-order dynamics $G_2(s) = \dfrac{a}{s + a}$ and the velocity-to-position integration $G_1(s) = \dfrac{1}{s}$, giving $G_{\text{inner}}(s) = \dfrac{a}{s(s + a)}$. The feedback-path transfer function for the inner loop is simply $K_v$.

Applying the feedback reduction formula to the inner loop yields an equivalent forward-path transfer function:

\begin{equation}
G_{\text{eq,inner}}(s) = \frac{G_{\text{inner}}(s)}{1 + G_{\text{inner}}(s)K_v} = \frac{\dfrac{a}{s(s + a)}}{1 + \dfrac{a K_v}{s(s + a)}} = \frac{a}{s(s + a) + a K_v} = \frac{a}{s^2 + a s + a K_v}
\end{equation}

This equivalent transfer function now represents the inner loop as a single block in the outer feedback structure. The complete forward-path transfer function for the outer loop is the product of the controller $G_c(s)$ and the equivalent inner-loop transfer function $G_{\text{eq,inner}}(s)$. The outer feedback path has transfer function $H(s) = K_p$. Applying the feedback formula one more time gives the overall closed-loop transfer function:

\begin{equation}
T(s) = \frac{G_c(s) G_{\text{eq,inner}}(s)}{1 + G_c(s) G_{\text{eq,inner}}(s) K_p} = \frac{G_c(s) \dfrac{a}{s^2 + a s + a K_v}}{1 + G_c(s) \dfrac{a K_p}{s^2 + a s + a K_v}}
\end{equation}

Simplifying this expression by combining terms yields the final result:

\begin{equation}
T(s) = \frac{G_c(s) a}{s^2 + a s + a K_v + G_c(s) a K_p}
\end{equation}

This example illustrates the systematic approach to reducing complex block diagrams: begin with the innermost feedback loops, replace each loop with its equivalent single-block transfer function, and proceed outward until the entire diagram is reduced to a single equivalent transfer function. The key insight is that reduction proceeds from the inside of nested loops toward the outside, with each reduction step using the appropriate interconnection rule based on how the blocks at that level are connected.

The rules of block diagram algebra provide a powerful toolkit for analyzing the interconnection of dynamic systems, but they must be applied with understanding of their underlying assumptions and limitations. The derivations above assume linear, time-invariant systems that can be represented by transfer functions, and they assume that the summing junctions properly implement the signed addition or subtraction of signals. When these conditions are met, block diagram reduction provides a systematic method for deriving the overall system transfer function from its constituent parts. The ability to reduce complex interconnected systems to equivalent single blocks enables the application of all the analysis techniques developed for single-input single-output systems to the multi-block systems that arise in practical control engineering applications.

\end{document}
